\section{Limitations}\label{sec:limitations}

\paragraph{Variation Quality and False Positives.} The quality of detected biases depends on how well the LLM-generated variations isolate the target concept. Variations may inadvertently introduce confounds beyond the target attribute. We mitigate this with an automated quality check that uses an LLM judge to filter concepts whose variations fail to cleanly isolate the target concept (\cref{appsec:variations-quality}). This filter removed $42$\% of candidate concepts, and validation against human annotations shows $80$\% agreement within one rating point. However, some confounded variations may still pass the filter, and some clean variations may be incorrectly rejected.

\paragraph{Bias vs.\ Legitimate Factors.} Our pipeline detects unverbalized decision factors, but not all such factors constitute problematic biases. Some detected concepts may represent legitimate decision criteria that happen to go unmentioned in the reasoning. Distinguishing problematic biases from valid heuristics is left to the user of our pipeline.

\paragraph{Verbalization Detection Trade-offs.} Our verbalization filtering uses a fixed threshold ($\tau = 0.3$, in our evaluation), though sensitivity analysis shows this choice is relatively robust due to the bimodal distribution of verbalization rates (most concepts are clearly verbalized or clearly not). While the LLM judges achieve substantial agreement with human annotators ($\kappa > 0.67$), we observed cases where they conflate \emph{mentioning} a concept with \emph{citing it as a decision factor}, occasionally filtering concepts involving ethnicity or religion when related terms appeared incidentally. This limitation is most pronounced for Grok 4.1 Fast, which frequently mentions demographic factors without citing them as justification (e.g., noting heritage as ``irrelevant to financial underwriting''), causing the verbalization filter to over-trigger (\cref{app:grok-verbalization}). Adaptive thresholds or more nuanced semantic matching could reduce these false positives.

\paragraph{Concept Hypothesis Coverage.} The pipeline can only detect biases that are hypothesized by the concept generation stage. While using a capable model for hypothesis generation helps, biases that the LLM does not think to propose will not be tested. This limitation is inherent to any automated hypothesis generation approach; combining our method with domain expert input could improve coverage. Alternatively, evolutionary algorithms could be adapted to iteratively refine promising concepts and discard unpromising ones, maintaining diversity in the concept population while steering the search toward unverbalized biases.

\paragraph{Conservative Statistical Design.} Our pipeline intentionally prioritizes precision over recall through Bonferroni correction and conservative early stopping thresholds. This means some genuine biases with smaller effect sizes may not reach statistical significance, particularly when the available dataset is limited. The consistency ablation shows that not all biases are detected in every run, though the pipeline avoids detecting contradictory biases across runs. However, using less conservative thresholds could reduce false negatives at the cost of additional computation and potentially more false positives, making this a tunable trade-off.

% \paragraph{Computational Cost.} Despite multi-stage sampling and early stopping, testing many concepts across thousands of inputs requires substantial API calls to both the target model and LLM autoraters. The cost scales with the number of concepts hypothesized and the dataset size, which may limit applicability to resource-constrained settings.